% GENERATED FILE. Edit canonical lessons, then run npm run generate:books.
% canonical-source-hash: fnv1a64:0173cb984b01316e
% canonical-lessons: FR-C09-janvier, FR-C09-fevrier, FR-C09-mars, FR-C09-avril, FR-C09-mai, FR-C09-juin, FR-C09-mois-practice-1

\chapter{The Months: Gods and Beginnings}
\label{ch:months-gods}
\hlchaptermodality{9}

\begin{chapteropening}
\textbf{By the end of this chapter:} \emph{I can name the first six months in French and say which Roman figure each one carries.}

Six months, and not one of them is a number. The calendar is a procession of Roman gods and festivals carried almost unchanged out of Latin -- Janus in the doorway of the year, the Februa that cleaned the city before it turned, and the war-god who already owns a weekday.
\end{chapteropening}

\section[janvier]{janvier — January}

\label{lesson:FR-C09-janvier}



\begin{quote}
The week was a map of the sky. The calendar is a \textbf{procession of Roman gods and emperors}, carried almost unchanged out of Latin — and it opens with the one god who faces both ways at once.
\end{quote}

\begin{sounds}
\begin{itemize}
\raggedright
  \item \texttt{j-zh} — the soft \textbf{j} of \emph{measure}, the one you met in \textbf{jeudi}.
  \item \texttt{nasal-an} — the \textbf{an} is nasal.
  \item Together: \textbf{janvier} = \emph{zhan-vee-AY}, three syllables, weight at the end.
\end{itemize}
\end{sounds}

\begin{cousinweb}
\textbf{janvier} is Latin \textbf{Januarius}, \textquotedblleft{}the month of \textbf{Janus}.\textquotedblright{}

Janus is the god of \textbf{doorways} — of gates, of thresholds, of any moment where one thing stops and another starts. He is carved with \textbf{two faces}, one looking back and one looking forward, because that is what a doorway is.

Put the first month of the year in his care and the name explains itself: the month that stands in the door, looking back at the year that closed and forward into the one opening. English \textbf{January} is the identical word; a \textbf{janitor} was once a doorkeeper, from the same \emph{janua}, \textquotedblleft{}door.\textquotedblright{}
\end{cousinweb}

\subsection*{Your turn}
Say these aloud:

\begin{itemize}
\raggedright
  \item \textquotedblleft{}janvier\textquotedblright{} — \emph{zhan-vee-AY}, the soft \emph{j} of \textbf{jeudi}
  \item the god — \textquotedblleft{}Janus, two faces, the doorway of the year\textquotedblright{}
  \item \textquotedblleft{}janvier\textquotedblright{} then English \textquotedblleft{}January, janitor\textquotedblright{} — one door-god in all three
\end{itemize}

\begin{tcolorbox}[breakable,colback=teal!4,colframe=teal!35!black,title={Before you move on}]
How do you say January? (\textbf{janvier}, \emph{zhan-vee-AY}.) Which god is it named for, and what is he god of? (\textbf{Janus} — \textbf{doorways}, beginnings.) Why does he have two faces? (One looks \textbf{back}, one looks \textbf{forward} — the year's threshold.) Next: the month named for a washing.
\end{tcolorbox}
\section[février]{février — February}

\label{lesson:FR-C09-fevrier}



\begin{quote}
Eleven of the twelve months are named for a god, an emperor or a number. This one is named for a \textbf{festival} — and specifically for a cleaning.
\end{quote}

\begin{sounds}
\begin{itemize}
\raggedright
  \item \texttt{vowel-e-acute} — the \textbf{é} is a tight, bright \emph{ay}, said with the mouth narrow. It is the accent that tells you the vowel is that one and no other.
  \item \textbf{février} = \emph{fay-vree-AY}. The \textbf{r} is the throat-\textbf{r}, and the ending is the same \emph{-AY} as \textbf{janvier}.
\end{itemize}
\end{sounds}

\begin{culture}
\textbf{février} is Latin \textbf{Februarius}, from the \textbf{Februa} — a Roman festival of \textbf{purification} held in this month.

That placement is the interesting part. The old Roman year began in \textbf{March}, so February was the \textbf{last} month, and the city cleaned itself before the year turned. The month is named for the scrubbing that came before the fresh start.

English \textbf{February} is the same word, and English speakers have been quietly dropping its first \textbf{r} for centuries.
\end{culture}

\subsection*{Your turn}
Say these aloud:

\begin{itemize}
\raggedright
  \item \textquotedblleft{}février\textquotedblright{} — \emph{fay-vree-AY}, the \emph{é} bright and narrow
  \item \textquotedblleft{}janvier, février\textquotedblright{} — the first two, both ending \emph{-AY}
  \item the reason — \textquotedblleft{}the Februa, the cleaning before the year turned\textquotedblright{}
\end{itemize}

\begin{tcolorbox}[breakable,colback=teal!4,colframe=teal!35!black,title={Before you move on}]
How do you say February? (\textbf{février}, \emph{fay-vree-AY}.) What is it named after? (The \textbf{Februa}, a festival of \textbf{purification}.) Why did a cleaning-festival fall at the end of the year? (The Roman year began in \textbf{March}, so February closed it.) Next: a month you already know as a day.
\end{tcolorbox}
\section[mars]{mars — March}

\label{lesson:FR-C09-mars}



\begin{quote}
This one you have already half-learned, in a different part of the calendar. The war-god who owns Tuesday also owns a month.
\end{quote}

\begin{sounds}
\begin{itemize}
\raggedright
  \item \textbf{mars} = \emph{marss}. The final \textbf{s} \textbf{is} sounded here — one of the small set of French words that keeps it.
  \item Do not read it as English \emph{Mars}, which has a different vowel; the French \textbf{a} is open and flat.
\end{itemize}
\end{sounds}

\begin{cousinweb}
\textbf{mars} is Latin \textbf{Martius}, the month of \textbf{Mars}, the war-god.

Now put it beside the weekday:

\begin{itemize}
\raggedright
  \item \textbf{mars} — \emph{Martius}, \textbf{his month}
  \item \textbf{mardi} — \emph{Martis diēs}, \textbf{his day}
\end{itemize}

\textbf{The same god, twice in the calendar.} Spot one and you already know the other, which is the first time a French word has paid you back in two different systems at once. English keeps the god's name in \textbf{martial}, \textquotedblleft{}of war,\textquotedblright{} and in the planet \textbf{Mars}.
\end{cousinweb}

\begin{culture}
March was the \textbf{first month} of the old Roman year. That single fact has now explained three separate things you have met: why \emph{février}'s cleaning festival came at the end, why the counting months land two places late, and why the war-god gets the opening of the year — campaigning season began when the weather turned.
\end{culture}

\subsection*{Your turn}
Say these aloud:

\begin{itemize}
\raggedright
  \item \textquotedblleft{}mars\textquotedblright{} — \emph{marss}, the \emph{s} sounded
  \item the pair — \textquotedblleft{}mars, mardi\textquotedblright{}: his month and his day
  \item \textquotedblleft{}janvier, février, mars\textquotedblright{} — the first three
\end{itemize}

\begin{tcolorbox}[breakable,colback=teal!4,colframe=teal!35!black,title={Before you move on}]
How do you say March? (\textbf{mars}, \emph{marss}, with the \emph{s}.) Which god is it named for? (\textbf{Mars}, the war-god.) Which weekday is named for the same god, and what is the difference between the two words? (\textbf{mardi} — \emph{Martis diēs}, his \textbf{day}; \emph{mars} is \emph{Martius}, his \textbf{month}.) Which month opened the old Roman year? (\textbf{March}.) Next: the month the buds open.
\end{tcolorbox}
\section[avril]{avril — April}

\label{lesson:FR-C09-avril}



\begin{quote}
Every month so far has had a firm answer behind it. This one has a guess — and the guess is old enough that the Romans were already making it.
\end{quote}

\begin{sounds}
\begin{itemize}
\raggedright
  \item \textbf{avril} = \emph{a-VREEL}. Two syllables; the \textbf{r} sits at the back of the throat, and the \textbf{i} is the tight vowel of \emph{machine}.
  \item The final \textbf{l} is sounded.
\end{itemize}
\end{sounds}

\begin{cousinweb}
\textbf{avril} is Latin \textbf{Aprilis}, and where \emph{Aprilis} came from is genuinely uncertain. The oldest guess, and still the best one, is Latin \textbf{aperīre}, \textquotedblleft{}to \textbf{open}\textquotedblright{} — the month the buds open.

It is worth noticing that this is a \textbf{guess}, and saying so. A book that hands you a confident story for every word is inventing some of them. English \textbf{aperture}, an opening, comes from that same \emph{aperīre}, so if the guess is right, an aperture and April are the same idea: something unclosing.
\end{cousinweb}

\subsection*{Your turn}
Say these aloud:

\begin{itemize}
\raggedright
  \item \textquotedblleft{}avril\textquotedblright{} — \emph{a-VREEL}, the \emph{l} sounded
  \item \textquotedblleft{}mars, avril\textquotedblright{} — the pair
  \item the guess — \textquotedblleft{}aperīre, to open; the buds open\textquotedblright{}
\end{itemize}

\begin{tcolorbox}[breakable,colback=teal!4,colframe=teal!35!black,title={Before you move on}]
How do you say April? (\textbf{avril}, \emph{a-VREEL}.) What is the best guess at its root, and what does it mean? (\textbf{aperīre}, \textquotedblleft{}to \textbf{open}.\textquotedblright{}) How certain is that? (Not very — it is a \textbf{guess}, and an ancient one.) Next: the goddess of growing.
\end{tcolorbox}
\section[mai]{mai — May}

\label{lesson:FR-C09-mai}



\begin{quote}
The shortest month-name in French, and one of the shortest words you have met: three letters, one syllable.
\end{quote}

\begin{sounds}
\begin{itemize}
\raggedright
  \item \texttt{ai-y} — the pair \textbf{ai} is a single sound, close to the \emph{e} of English \emph{bed}. So \textbf{mai} = \emph{may}, in one beat. Neither letter is heard on its own.
  \item This is the same \textbf{ai} you will meet again and again; it is one of the most common vowel pairs in French, and it never splits.
\end{itemize}
\end{sounds}

\begin{cousinweb}
\textbf{mai} is Latin \textbf{Maius}, the month of \textbf{Maia} — an old Italic goddess of \textbf{growth} and increase, worshipped long before the Romans wrote much down about her.

Her name is thought to sit on the same root as Latin \textbf{maior}, \textquotedblleft{}greater,\textquotedblright{} which English carries in \textbf{major}, \textbf{majority} and \textbf{mayor}. If that is right, then \emph{mai} is the \textbf{growing} month in the plainest sense: the one where things get bigger.
\end{cousinweb}

\subsection*{Your turn}
Say these aloud:

\begin{itemize}
\raggedright
  \item \textquotedblleft{}mai\textquotedblright{} — \emph{may}, one beat, the \emph{ai} never split
  \item \textquotedblleft{}avril, mai\textquotedblright{} — the pair
  \item the goddess — \textquotedblleft{}Maia, growth; and maior, greater\textquotedblright{}
\end{itemize}

\begin{tcolorbox}[breakable,colback=teal!4,colframe=teal!35!black,title={Before you move on}]
How do you say May? (\textbf{mai}, \emph{may}.) What sound does the pair \textbf{ai} make? (One sound, close to the \emph{e} of \emph{bed} — never two.) Which goddess is the month named for, and what is she goddess of? (\textbf{Maia} — \textbf{growth}.) Next: the queen of the gods.
\end{tcolorbox}
\section[juin]{juin — June}

\label{lesson:FR-C09-juin}



\begin{quote}
The sixth month closes the half-year, and it belongs to the one figure in the calendar who outranks everybody else in it.
\end{quote}

\begin{sounds}
\begin{itemize}
\raggedright
  \item \texttt{j-zh} — the soft \textbf{j} again, as in \textbf{jeudi} and \textbf{janvier}.
  \item \texttt{nasal-in} — the \textbf{in} is nasal, as in \textbf{cinq}.
  \item Together: \textbf{juin} = \emph{zhwa(n)}, one syllable. It is shorter than it looks.
\end{itemize}
\end{sounds}

\begin{cousinweb}
\textbf{juin} is Latin \textbf{Junius}, the month of \textbf{Juno} — queen of the gods, wife of Jupiter, and the goddess of marriage.

That last detail still shows. June is the traditional wedding month across Europe because it was \textbf{her} month, and couples wanted her on their side.

You have already met her husband: \textbf{jeudi}, \emph{Jovis diēs}, is Jupiter's day. So the calendar now holds both halves of the divine marriage — his \textbf{day} and her \textbf{month}.
\end{cousinweb}

\subsection*{Your turn}
Say these aloud:

\begin{itemize}
\raggedright
  \item \textquotedblleft{}juin\textquotedblright{} — \emph{zhwa(n)}, one syllable
  \item \textquotedblleft{}mai, juin\textquotedblright{} — the pair
  \item the couple — \textquotedblleft{}juin is Juno's month, jeudi is Jupiter's day\textquotedblright{}
\end{itemize}

\begin{tcolorbox}[breakable,colback=teal!4,colframe=teal!35!black,title={Before you move on}]
How do you say June? (\textbf{juin}, \emph{zhwa(n)}, one syllable.) Which goddess is it named for? (\textbf{Juno}, queen of the gods.) Why are weddings held in June? (She is the goddess of \textbf{marriage}, and it is her month.) Which day belongs to her husband? (\textbf{jeudi} — Jupiter's.) Next: the first six, run together.
\end{tcolorbox}
\section[Practice]{Practice — the first half of the year}

\label{lesson:FR-C09-mois-practice-1}



\begin{quote}
Six months, taken one at a time. Run them together now, then say what is inside each one.
\end{quote}

\subsection*{Your turn: the six, in order}
Read down the left, then cover the right and read back:

\noindent
\begin{tabularx}{\linewidth}{@{}>{\raggedright\arraybackslash}X>{\raggedright\arraybackslash}X@{}}
\toprule
\textbf{French} & \textbf{English} \\
\midrule
\textbf{janvier} & January \\
\textbf{février} & February \\
\textbf{mars} & March \\
\textbf{avril} & April \\
\textbf{mai} & May \\
\textbf{juin} & June \\
\bottomrule
\end{tabularx}

\emph{Twice through:} the six months, in order

Four of them end in that same unstressed \emph{-AY}; \textbf{mars} and \textbf{mai} are the two short ones, and they are the two easiest to lose.

\subsection*{Your turn: name what is inside}
Say each month, then the figure or thing behind it:

\begin{itemize}
\raggedright
  \item \textbf{janvier} — \textbf{Janus}, the two-faced god of doorways
  \item \textbf{février} — the \textbf{Februa}, a festival of purification
  \item \textbf{mars} — \textbf{Mars}, the war-god (and \emph{mardi} is his day)
  \item \textbf{avril} — perhaps \emph{aperīre}, \textquotedblleft{}to open,\textquotedblright{} though nobody is certain
  \item \textbf{mai} — \textbf{Maia}, goddess of growth
  \item \textbf{juin} — \textbf{Juno}, queen of the gods
\end{itemize}

\begin{grammarlens}[title={What you've built this chapter}]
\begin{itemize}
\raggedright
  \item \textbf{six month names}, each of which is a Roman figure rather than a number.
  \item \textbf{one word paying twice} — \emph{mars} and \emph{mardi} are the same god's month and day, which is the first time a French word has cashed out in two systems.
  \item \textbf{the old Roman year began in March}, which has now explained three separate oddities: why a cleaning-festival names February, why the war-god opens the year, and why four months you have seen are counting two places out.
\end{itemize}
\end{grammarlens}

\begin{tcolorbox}[breakable,colback=teal!4,colframe=teal!35!black,title={Before you move on}]
Give the six months without stopping. Now start from \textbf{avril}. Which month shares its god with a weekday, and which weekday? (\textbf{mars} — with \textbf{mardi}, both Mars.) Which month is named for a festival rather than a person? (\textbf{février}, the Februa.) Which one is a guess? (\textbf{avril} — perhaps \emph{aperīre}, \textquotedblleft{}to open.\textquotedblright{})
\end{tcolorbox}
